\documentclass[UTF8,fontset=fandol,11pt,a4paper,landscape]{ctexart}
\usepackage[margin=18mm,headheight=16pt]{geometry}
\usepackage{fontspec}
\setmainfont{texgyretermes-regular.otf}[BoldFont=texgyretermes-bold.otf]
\usepackage{fancyhdr,hyperref}
% Include in the document preamble. Native vector figures; no external artwork.
\usepackage{tikz}
\usetikzlibrary{arrows.meta,shapes.geometric,calc}
\definecolor{FlowBlue}{HTML}{245B78}
\definecolor{FlowLight}{HTML}{EDF5FA}
\definecolor{FlowAmber}{HTML}{A46B1C}
\definecolor{FlowWarm}{HTML}{FBF2E3}
\definecolor{FlowGray}{HTML}{69737A}
\tikzset{
  operation/.style={draw=FlowBlue,fill=FlowLight,rounded corners=3mm,
    minimum height=1.55cm,text width=3.35cm,align=center,
    font=\small,inner sep=2.5mm,line width=.7pt},
  focus/.style={operation,draw=FlowAmber,fill=FlowWarm},
  candidate/.style={operation,draw=FlowGray,fill=gray!6,dashed},
  branch/.style={diamond,aspect=1.8,draw=FlowBlue,fill=white,
    text width=1.7cm,align=center,font=\small,inner sep=1mm,line width=.7pt},
  flow/.style={-{Stealth[length=2.5mm]},draw=FlowGray,line width=.85pt},
  conditional/.style={flow,dashed},
  edgenote/.style={font=\footnotesize,fill=white,inner sep=1.2mm,align=center},
  statetag/.style={font=\footnotesize,text=FlowGray,align=center}
}

\hypersetup{pdftitle={可复用工艺图模板},pdfauthor={},colorlinks=true,linkcolor=FlowBlue}
\pagestyle{fancy}\fancyhf{}
\fancyhead[L]{\small\color{FlowGray}制造工艺与成本管理 · 图形模板 1.0}
\fancyhead[R]{\small\color{FlowGray}合成示例 · 按项目填写与复核}
\fancyfoot[L]{\small\color{FlowGray}节点和连线是可编辑占位；图形不代表工艺已经采用或能力已经验证。}
\fancyfoot[R]{\thepage}
\setlength{\parindent}{0pt}\setlength{\parskip}{6pt}
\newcommand{\FlowTitle}[2]{{\LARGE\bfseries\color{FlowBlue}#1\par}\vspace{2mm}{\color{FlowGray}#2\par}\vspace{7mm}}
\newcommand{\FlowNotes}[1]{\vfill\small\textbf{填写与复核：}#1\par\vspace{2mm}\textcolor{FlowGray}{图例：蓝框为一般节点；琥珀框为本图强调节点；灰色虚框为候选或待定项；菱形为条件判断。颜色本身不表示质量等级或批准状态。}}
\begin{document}
\FlowTitle{01\quad 顺序工序与产品状态}{用途：表达一条已选路线的工序先后，并明确检验面对哪一种状态。}
\begin{center}\begin{tikzpicture}[x=1cm,y=1cm]
  \node[operation] (a) at (0,0) {OP10\quad 工序 A\\输入 S0；输出 S1};
  \node[operation] (b) at (5,0) {OP20\quad 工序 B\\输入 S1；输出 S2};
  \node[operation] (c) at (10,0) {OP30\quad 工序 C\\输入 S2；输出 S3};
  \node[focus] (i) at (15,0) {IN40\quad 阶段检验\\受检状态 S3\\方法与判据待填};
  \node[operation] (n) at (20,0) {OP50\quad 下一阶段\\进入条件待填};
  \draw[flow] (a)--(b);
  \draw[flow] (b)--(c);
  \draw[flow] (c)--(i);
  \draw[flow] (i)--node[edgenote,above=11mm]{按放行范围}(n);
  \node[statetag] at (10,-2.1) {工序编号用于追溯；箭头表达本路线的状态先后关系。};
\end{tikzpicture}
\end{center}
\FlowNotes{替换工序 A/B/C 及状态 S0 至 S3；登记对象版次、来源和检验阶段。后续工序改变已检特性时，重新核查检验覆盖。具体清洗、热处理或表面处理位置依项目要求填写。}
\clearpage
\FlowTitle{02\quad 并行加工、汇合与装配}{用途：表达多个对象分别投入，以及进入汇合工序必须满足的条件。}
\begin{center}\begin{tikzpicture}[x=1cm,y=1cm]
  \node[operation] (a) at (0,2) {OP10A\quad 零件 A 加工\\输出 A1};
  \node[operation] (b) at (0,-2) {OP10B\quad 零件 B 加工\\输出 B1};
  \node[focus] (j) at (5,0) {OP20\quad 汇合准备\\核对身份、版次与接口};
  \node[operation] (c) at (10,0) {OP30\quad 装配或连接\\A1＋B1 → AB1};
  \node[focus] (i) at (15,0) {IN40\quad 阶段检验\\受检对象 AB1};
  \node[operation] (n) at (20,0) {OP50\quad 下道工序\\进入条件待填};
  \draw[flow] (a.east)--(2.7,2)--(2.7,0)--(j.west);
  \draw[flow] (b.east)--(2.7,-2)--(2.7,0)--(j.west);
  \draw[flow] (j)--node[edgenote,above=10mm]{输入均到位}(c);
  \draw[flow] (c)--(i);
  \draw[flow] (i)--(n);
  \node[statetag] at (10,-3.1) {并行投入不代表加工节拍相同；汇合条件和对象状态分别登记。};
\end{tikzpicture}
\end{center}
\FlowNotes{分别标识零件 A/B 的版次和输出状态；填写两个输入是否都必需、接口检查内容及装配后的新对象。根据实际路线替换“装配或连接”，不要由图形相似直接移植工艺。}
\clearpage
\FlowTitle{03\quad 检验、不合格处置与返工复验}{用途：把正常路径、未决处置和获准后的返工路径同时表达清楚。}
\begin{center}\begin{tikzpicture}[x=1cm,y=1cm]
  \node[focus] (i) at (0,0) {IN10\quad 阶段检验\\对象、状态、判据待填};
  \node[branch] (d) at (5,0) {符合\\判据？};
  \node[operation] (n) at (10,0) {OP20\quad 下一阶段\\按规定范围放行};
  \node[operation] (h) at (5,-3) {NC10\quad 隔离与评审\\处置尚未确定};
  \node[branch] (a) at (10,-3) {返工获\\批准？};
  \node[candidate] (o) at (15,-3) {NC20\quad 其他处置\\另行决定与记录};
  \node[operation] (r) at (10,-6) {RW10\quad 按批准路线返工\\对象与路线版次待填};
  \node[focus] (t) at (0,-6) {IN30\quad 受影响特性复验\\范围与判据待填};
  \draw[flow] (i)--(d);
  \draw[flow] (d)--node[edgenote,above]{是}(n);
  \draw[flow] (d)--node[edgenote,right]{否}(h);
  \draw[flow] (h)--(a);
  \draw[flow] (a)--node[edgenote,above=10mm]{否／待定}(o);
  \draw[flow] (a)--node[edgenote,right]{是}(r);
  \draw[flow] (r)--node[edgenote,above]{返工完成}(t);
  \draw[flow] (t.west)--(-2.7,-6)--(-2.7,2)--node[edgenote,above]{复验结果}(5,2)--(d.north);
\end{tikzpicture}
\end{center}
\FlowNotes{填写受检对象、状态、阶段、方法与判据。返工须有批准路线；其他处置另行决定。复验范围按受影响特性与既定要求确定，不能把返工完成直接画成检验通过。}
\clearpage
\FlowTitle{04\quad 可选工序与未决分支}{用途：资料尚不齐全时保留条件路线，让已采用、候选和待澄清保持可见。}
\begin{center}\begin{tikzpicture}[x=1cm,y=1cm]
  \node[operation] (s) at (0,0) {OP10\quad 当前工序完成\\输出状态 S1};
  \node[branch] (d) at (5,0) {附加工序\\是否适用？};
  \node[candidate] (p) at (5,3.3) {OP20\quad 候选附加工序\\采用依据待确认};
  \node[focus] (i) at (10,3.3) {IN30\quad 受影响特性检查\\受检状态、判据待填};
  \node[operation] (st) at (10,0) {ST40\quad 登记实际输出状态\\关联已采用路线};
  \node[operation] (n) at (15,0) {OP50\quad 后续工序\\核对进入条件};
  \node[candidate] (q) at (5,-3) {Q10\quad 澄清或验证\\记录未决条件};
  \draw[flow] (s)--(d);
  \draw[conditional] (d.north)--node[edgenote,right]{适用且已采用}(p.south);
  \draw[conditional] (p)--(i);
  \draw[flow] (i)--(st);
  \draw[flow] (d)--node[edgenote,above=1mm,text width=10mm,inner sep=.5mm]{确认\\不适用}(st);
  \draw[flow] (st)--(n);
  \draw[conditional,<->] (d.south)--node[edgenote,right]{待确认／补充依据}(q.north);
\end{tikzpicture}
\end{center}
\FlowNotes{填写附加工序的适用条件、来源和采用记录；未确认不能自行按“不适用”执行。附加工序改变状态或特性时，检查对应的检验覆盖。虚线位置不能代替制造顺序或批准记录。}
\clearpage
\FlowTitle{05\quad 工序、资源、证据与成本关系}{用途：解释方案依据和费用从哪里来。\textbf{本图箭头表示知识关系，不表示制造顺序。}}
\begin{center}\begin{tikzpicture}[x=1cm,y=1cm]
  \node[operation] (p) at (0,2) {OBJ-01\quad 产品对象\\版次、状态与要求};
  \node[operation] (op) at (6,2) {OP-10\quad 工序或检验\\输入、输出与判据};
  \node[focus] (c) at (12,2) {C-10\quad 费用项\\按件／按批／一次性};
  \node[operation] (v) at (18,2) {VIEW-01\quad 条件预算\\数量、币税与范围};
  \node[operation] (e) at (0,-2) {SRC-01\quad 证据来源\\原件位置与适用条件};
  \node[candidate] (r) at (6,-2) {RES-01\quad 资源候选\\设备、工装或供应方};
  \node[operation] (f) at (12,-2) {FB-01\quad 实际反馈\\节拍、返工与费用};
  \draw[flow] (p)--node[edgenote,above]{提出需求}(op);
  \draw[flow] (op)--node[edgenote,above]{发生费用}(c);
  \draw[flow] (c)--node[edgenote,above]{纳入核算}(v);
  \draw[flow] (e)--node[edgenote,left]{支持记录}(p);
  \draw[flow] (e)--node[edgenote,above=10mm]{支持能力声明}(r);
  \draw[conditional] (r)--node[edgenote,left,align=right]{承担工序\\待适配验证}(op);
  \draw[flow] (op.south east)--node[edgenote,sloped,above]{执行后回填}(f.north west);
  \draw[flow] (f)--node[edgenote,right]{修订估算}(c);
  \node[statetag] at (9,-3.7) {本图箭头为知识关系，不表示制造顺序。};
\end{tikzpicture}
\end{center}
\FlowNotes{把节点标识对应到本项目的对象、工序、资源、费用和来源。资源可调配与工序能力分开记录；准备费是否已含、数量口径和实测范围分别填写。实际反馈支持修订估算，不自动证明已经降本。}
\end{document}
